\title{Large Language Models Can Automatically Engineer Features for Few-Shot Tabular Learning}

\begin{abstract}
Large Language Models (LLMs) have recently demonstrated exceptional proficiency in solving novel and complex reasoning problems, making them a promising avenue for advancing tabular learning—a task essential to numerous practical applications. In this work, we introduce a new in-context learning framework termed \textsf{FeatLLM}, which leverages LLMs to serve as automated feature engineers. This framework generates an input feature set tailored for tabular prediction tasks. The features synthesized through this method are subsequently utilized by a basic downstream machine learning model, such as linear regression, thereby facilitating effective few-shot learning with high predictive accuracy. Importantly, during inference, the \textsf{FeatLLM} approach relies solely on this simple predictive model, negating the need for recurrent LLM queries. In contrast to other LLM-centric methods, our approach does not require LLM invocation for each test instance, thereby reducing computational overhead. It also operates with basic API access and bypasses constraints linked to maximum prompt lengths. Through extensive experiments across various tabular datasets spanning multiple fields, we show that \textsf{FeatLLM} reliably discovers high-quality rules, yielding considerable (on average 10\%) performance gains relative to alternatives like TabLLM and STUNT.
\end{abstract}

\section{Introduction}
Large Language Models (LLMs) have attained remarkable success in producing appropriate outputs for previously unseen tasks without the necessity for domain-specific fine-tuning~\cite{creswell2022selection,imani2023mathprompter,taylor2022galactica}. Trained on vast collections of textual data, these models encapsulate substantial prior knowledge, effectively standing in for specialized expertise across numerous fields~\cite{Genomes2021,singhal2023large,wilcox2003role}. This characteristic has proven essential for streamlining a range of applications, including dialogue analysis~\cite{finch2023leveraging} and automated program repair~\cite{fan2023automated}. More specifically, by formulating prompts that combine the task description with several illustrative examples, LLMs are capable of discerning contextual information, prioritizing salient features and relevant instances~\cite{brown2020language,zhao2023survey}. Such observations highlight the strong analytical capacity of LLMs, pointing to their significant potential in the domain of automated feature engineering.

Recent developments have extended the exploitation of LLMs’ prior knowledge and reasoning to tasks within tabular learning. For example, domain knowledge—such as correlations between advanced age and higher disease risk—can inform models in medical prediction scenarios. Contemporary methodologies have translated task specifications and feature information into natural language, serialized the resulting data, and subsequently provided these prompts to LLMs~\cite{dinh2022lift,wang2023anypredict}. Following this, strategies such as in-context learning~\cite{nam2023semi} or more efficient model fine-tuning techniques~\cite{dinh2022lift,hegselmann2023tabllm} are typically employed. By harnessing the implicit expertise encoded within LLMs, these methods have surpassed traditional tabular learning techniques—especially in data-scarce settings~\cite{hegselmann2023tabllm}.

Existing methods that leverage LLMs for tabular learning face several inherent drawbacks. End-to-end prediction pipelines typically demand at least a single LLM invocation per datapoint, which translates to high computational costs. Moreover, most strategies necessitate fine-tuning the LLM to achieve optimal performance, thereby limiting their applicability when training frameworks or tuning interfaces are unavailable—an issue particularly salient with newly-developed, state-of-the-art LLMs that are only accessible through restricted inference APIs~\cite{anil2023palm,openai2023gpt4}. Additionally, when faced with lengthy prompts—such as those arising from high-dimensional tabular datasets—most current techniques become impractical or experience diminished accuracy~\cite{liu2023lost}. These limitations hinder the deployment of LLM-based tabular learning solutions in practical settings.

To address these concerns, we propose a paradigm shift in LLM utilization: rather than employing LLMs as direct predictors, we harness their capacity for feature engineering. Our new approach, embodied in the framework \textsf{FeatLLM}, leverages in-context learning to decipher the "criteria" underlying LLM decision-making. For example, in a clinical prediction scenario, an LLM may generate explicit rules linking feature patterns to the presence of disease. These generated criteria are subsequently used to construct alternative features, supplanting the original inputs for downstream tasks. This reorientation of the LLM's role facilitates a more streamlined downstream process; once features are synthesized, the need for sophisticated prediction models is greatly diminished, thus improving inference speed. By exploiting the LLMs' proficiency in in-context learning, we can derive features without conventional training, merely by augmenting the prompt with illustrative examples. Finally, adopting ensemble-inspired strategies such as feature bagging~\cite{breiman2001random} enables us to circumvent the problem of overly long prompts.

\textsf{FeatLLM} organizes the feature generation prompt into two principal sub-tasks: (1) interpreting the problem and identifying how features relate to targets, and (2) leveraging insights gained and using limited-shot exemplars to formulate discriminative rules for class separation. This deliberate reasoning framework enables large language models to disregard irrelevant features during rule derivation. The resulting rules are subsequently implemented within the dataset (translating them into executable code), thus converting the data into binary indicators for each sample’s compliance with a rule. A low-complexity model is then trained to predict class probabilities using these engineered features. Robustness is enhanced through an ensemble strategy, in which feature generation and feature bagging are repeatedly performed. By tuning the count of features and samples selected for bagging, issues arising from excessive prompt length are effectively managed. The method operates without necessitating training phases and incurs minimal inference overhead, making LLM deployment practical in numerous real-world scenarios.

We assess the effectiveness of \textsf{FeatLLM} across 13 distinct tabular datasets under low-data settings, where it demonstrates reliable and superior performance. The results indicate that LLMs adeptly incorporate both prior knowledge and patterns gleaned from the data, successfully constructing valuable features. Our method consistently surpasses leading few-shot learning baselines under multiple conditions. The codebase is publicly available via an anonymized GitHub repository at~\texttt{https://github.com/Sungwon-Han/FeatLLM}.

\section{Related Work}

\subsection{Few-Shot Learning with Tabular Data}

Recent progress in few-shot learning has facilitated the extraction of broadly applicable representations from datasets while keeping annotation requirements low~\cite{chen2018closer}. While the technique originated with image analysis, its effectiveness now spans various modalities, encompassing not only text and audio but, more recently, tabular datasets as well~\cite{majumder2022few,nam2023stunt,schick2021exploiting}. A major challenge of learning from sparse labeled examples is the danger of picking up on non-generalizable, spurious patterns~\cite{bartlett2021deep}. To overcome these pitfalls, contemporary research has integrated auxiliary information or domain-relevant priors, shaping model learning via suitable inductive biases. Approaches include capitalizing on vast, unlabeled tabular datasets through sophisticated data augmentation for semi-supervised training~\cite{ucar2021subtab,yoon2020vime}, or leveraging unsupervised, task-agnostic pre-training to achieve better initializations~\cite{somepalli2021saint}. Such strategies frequently employ objectives that involve deliberate data corruption or masking with subsequent reconstruction~\cite{arik2021tabnet,majmundar2022met}, or utilize data augmentation for contrastive learning~\cite{bahri2022scarf,chen2020simple}. However, the diversity and complexity inherent in tabular formats complicate the discovery of consistent, general-purpose augmentations, and thus these techniques have shown limited effectiveness when faced with very limited labeled data~\cite{nam2023stunt}.

Another promising direction involves pre-training on a diverse set of tabular datasets—including those unrelated to the final application—and then transferring knowledge to downstream target tasks~\cite{zhang2023generative,levin2022transfer,zhu2023xtab}. For example, the TransTab model~\cite{wang2022transtab} employs transformer architectures that explicitly capture semantic associations between column names in addition to tabular cell values, enabling transfer of knowledge about tabular structure and content across tables. As a result, these models can perform few-shot or even zero-shot inference on novel tasks involving previously unseen tables. Similarly, TabPFN~\cite{hollmann2022tabpfn} synthesizes new tabular datasets from a blend of distributions for pre-training a Bayesian neural network. At inference, both training and test data from the downstream task are given as inputs to the model, which requires no further fine-tuning. Exposure to numerous tables and distributions equips these pre-trained models with inductive priors that outperform classical tree-based methods, especially in few-shot conditions.

\subsection{Language-Interfaced Tabular Learning}

Integrating large language models (LLMs)~\cite{anil2023palm,openai2023gpt4} represents a promising strategy for harnessing prior knowledge in tabular learning tasks. Trained on massive and varied text data, LLMs exhibit robust generalization abilities to new, previously unseen problems, and have found widespread success in diverse domains~\cite{brown2020language}. Recent research has sought to apply the latent knowledge contained within LLMs to tabular settings by encoding structured data as text strings that can be presented as prompts to these models~\cite{dinh2022lift,hegselmann2023tabllm,wang2023anypredict}. These prompts often amalgamate the actual tabular entries, metadata such as feature descriptions, and explicit task requirements, enabling the model to infer relevant relationships among inputs for prediction or reasoning tasks. The research landscape encompasses parameter-efficient adaptation methods, such as fine-tuning selected components with algorithms like LoRA~\cite{hu2021lora} or IA3~\cite{liu2022few}, as well as in-context learning paradigms wherein the model, without any update to its parameters, is prompted using a few labeled instances as demonstration—bypassing additional training~\cite{dinh2022lift,hegselmann2023tabllm,nam2023semi,slack2023tablet}.

Although the methods outlined above demonstrate considerable success in enhancing few-shot learning with tabular data, relying on LLMs for inference at all times introduces limitations, particularly for practical use within industrial environments. This research seeks to shift away from the typical black-box paradigm in which each prediction is processed entirely by an LLM. Rather than evaluating individual samples through the LLM, the approach focuses the LLM on deducing class-specific rules or conditions. \textsf{FeatLLM} translates these generated rules into programmatic code, producing new features that enable predictions to be made without repeated LLM inference. The method leverages in-context learning, using available prior knowledge and information from few-shot instances to derive relevant rules.

\section{Methods}
\textsf{FeatLLM} is built to identify and extract ``rules," corresponding to conditions for each possible class label, based on limited input samples, as illustrated in Figure~\ref{fig:main_model}. The system utilizes an LLM to generate such rules, which are subsequently converted into binary features for each sample—indicating compliance with each rule. These binary features are then subjected to a dot product with non-negative weights, which are adjusted through training to reflect each feature's contribution to class probability (see Section~\ref{Sec:training}). This workflow is executed multiple times using bagging, followed by aggregation through ensemble techniques. The subsequent sections outline the precise problem formulation and elaborate on each methodological component.

\vspace*{-0.12in}
\paragraph{Problem formulation.} We examine a tabular dataset containing $N$ labeled entries, $\mathcal{D} = \{(\mathbf{x}^i, \mathbf{y}^i)\}_{i=1}^{N}$, where each data point $\mathbf{x}^i$ is a vector in $d$ dimensions. The dataset features $d$ input variables, each uniquely named in natural language as $F = \{f_j\}_{j=1}^{d}$, with corresponding feature definitions provided elsewhere. Under the classification paradigm, every label $\mathbf{y}^i$ assigns the sample to one class among several predetermined categories. For $k$-shot learning, only $k$ samples, where $k$ is less than $N$, are selected at random for model training. 

\subsection{Prompt Design for Extracting Rules}
\label{Sec:prompt_design}
To support LLMs in deriving rules with more precise logical steps, we steer the process to reflect strategies an expert would adopt in tabular data analysis. Our prompting method encompasses three central components, illustrated in Fig.~\ref{fig:prompt_for_rule} and further detailed in Appendix~\ref{sec:example_rule}:

\vspace*{-0.12in}
\paragraph{Basic information description.} This segment supplies crucial details required to address the task (indicated in orange in Fig.~\ref{fig:prompt_for_rule}). Included are the articulation of the classification problem, label information, feature summaries, and representative examples. The task itself is introduced via a question format (for instance: ``Does this patient's myocardial infarction complications data suggest chronic heart failure? Yes or no?"). Additional task descriptions can be found in Appendix~\ref{sec:task_desc}. Feature descriptions specify the datatype for each variable and may optionally elaborate with the feature's definition. For categorical features, example category values can be listed. Demonstrative training samples are provided as text, annotated with their correct labels, following a serialization structure as used in prior work~\cite{dinh2022lift,hegselmann2023tabllm}:

\begin{align}
&\text{Serialize}(\mathbf{x}^i,\mathbf{y}^i, F) = \nonumber \\ 
&``\ f_1\ \text{is}\ \mathbf{x}^i_1.\ ... \ f_d\  \text{is}\  \mathbf{x}^i_d.\ \ \text{Answer:}\  \mathbf{y}^i\ ”
\end{align}

\vspace*{-0.12in}
\paragraph{Reasoning instruction.} Our approach seeks to bolster the reasoning abilities of the LLM by supplying explicit reasoning instructions (refer to blue text in Fig.~\ref{fig:prompt_for_rule}). The reasoning instruction starts with an opening sentence in the style of chain-of-thought prompting~\cite{wei2022chain}. Subsequently, the LLM's reasoning procedure is partitioned into two distinct stages. During the initial stage, the LLM is prompted to deduce the causal connections or general tendencies between features and the overall task description. This stage is completed independently from any example demonstrations, utilizing only its internal knowledge base and common sense to draw relevant prior insights about the problem at hand. In the following stage, the LLM then incorporates example demonstrations, in conjunction with the information extracted previously, to derive classification rules tailored to each class. This bifurcated method mitigates the risk of the model learning misleading correlations in unrelated columns and facilitates focusing on the most pertinent features. To strike a balance and avoid extracting an insufficient or overwhelming number of rules—which could either hinder model performance or result in an excessively large prompt—we require a predetermined number of rules (specifically, 10)\footnote{An in-depth discussion on the choice of rule count is provided in Section~\ref{sec:analysis}.} for every class. Moreover, the LLM is instructed to consult the feature type descriptions found in the basic information section when generating rules to ensure compatibility and to prevent mismatches in rule construction.

\vspace*{-0.12in}
\paragraph{Response instruction.} For efficient interpretation and utilization of the generated rules, the LLM receives dedicated instructions concerning response structuring (see yellow text in Fig.~\ref{fig:prompt_for_rule}). The prompt outlines precise directives regarding the appropriate formatting of responses for each answer class (i.e., response format), as well as the required organization of individual rules (i.e., rule format). These formatting guidelines enable the LLM to select suitable forms contingent on the feature types involved. Even without additional prompts, the LLM is capable of combining several rules via logical connectors such as AND and OR. An illustrative outcome for this can be found in Appendix~\ref{sec:outcome-rule}.

\subsection{Inferring Class Likelihood via Rules}
\label{Sec:training}

\paragraph{Parsing rules for feature generation.} In this stage, the rules obtained previously are employed to construct new binary features. Each feature corresponds to a specific class and indicates whether a sample satisfies the relevant class rules, outputting either '0' or '1'. These features are subsequently used to estimate the likelihood of a sample belonging to each respective class. Since the rules derived from the LLM are expressed in natural language, an automated parsing method is necessary for feature extraction. Although we attempt to enforce uniform formatting of rules and responses during their generation to facilitate easier parsing, the LLM’s output can still exhibit inconsistent syntactic modifications~\cite{kovavc2023large}. For example, the LLM may rephrase symbols, such as replacing ``$>$” and ``$<$” with ``greater than" or ``less than," thereby complicating the parsing process.

To mitigate parsing difficulties introduced by textual noise, we opt to use the LLM as a parsing tool rather than devising elaborate programmatic parsers. The LLM excels at capturing semantic meaning despite syntactic variability and can produce concise code based on provided instructions (benefiting from its exposure to code during training). To utilize these capabilities, we embed information such as the function name, descriptions of inputs and outputs, and the relevant rules within the prompt that is passed to the LLM. The prompt templates and sample outcomes are illustrated in Figures~\ref{fig:prompt_for_parsing} and~\ref{sec:example_parsing}-\ref{sec:example_parse_outcome} in the Appendix. The generated code snippets are executed via Python’s \texttt{exec()} function to accomplish the required data transformation.

\vspace*{-0.12in}
\paragraph{Inferring class likelihood.} After generating class-specific binary features from the rules, one straightforward way to estimate the likelihood that a sample belongs to a certain class is by summing the satisfied rules for each class—that is, accumulating the corresponding binary values per class. Nevertheless, the impact of individual rules and features is not uniform; their weights must be learned from labeled training data. To ascertain the relevance of each rule, we train a conventional machine learning (ML) model on the binary features for each class. A linear model without a bias term, for instance, can be utilized. Specifically, for sample $\mathbf{x}^i$ and class $k$, let $\mathbf{z}_k^i \in \{0, 1\}^R$ denote the vector of binary features (where $R$ is the count of class-specific rules and $c$ represents the number of classes). The predicted probability $\mathbf{p}^i$ for each class is then obtained as follows:

\begin{align}
\text{logit}_k^i &= \max(\mathbf{w}_k, 0) \cdot \mathbf{z}_k^i, \nonumber \\
\mathbf{p}^i &= \text{Softmax}({[\text{logit}_{1}^i, ..., \text{logit}_{c}^i]}). \label{eq:infer_prob}
\end{align}

In this context, $\mathbf{w}_k$ denotes the adjustable parameters within the linear model, representing the relative significance of each feature. By constraining these weights to a positive domain, it is possible to guarantee that the features constructed by the LLM do not diminish the probability associated with any class during the training phase. Subsequently, the logits are transformed into class probabilities by applying the softmax function. The cross-entropy loss serves as the optimization objective for tuning the weights $\mathbf{w}_k$, using the limited available few-shot training data.

\vspace*{-0.12in}
\paragraph{Ensembling with bagging.} To finalize predictions, the described procedure is iteratively performed to produce multiple inference models, whose outputs are then combined through ensemble techniques. With a set of trained weights from $T$ independent runs, denoted as $\{\mathbf{w}[t]\}_{t=1}^T$, the classification outcome is computed by taking the mean of the per-class probabilities $\mathbf{p}^i$ obtained from each weight instance $\mathbf{w}[t]$ in accordance with Eq.~\ref{eq:infer_prob}.

To foster variation in rule generation for the ensemble process, the temperature parameter for the LLM inference is set sufficiently high, promoting stochasticity. Additionally, the sequence in which few-shot demonstrations are presented is randomized, based on empirical evidence indicating that their order can alter LLM responses~\cite{lu2022fantastically}\footnote{Appendix~\ref{sec:diversity} presents a detailed analysis of rule diversity}. To leverage prediction variability for increased resilience, bagging is employed, selecting different subsets of features or samples in each trial—a strategy that is especially advantageous with a larger training corpus or feature set. The advocated ensemble methodology delivers two primary benefits. Firstly, errors arising from incorrectly derived rules in certain trials can be mitigated by correct rules from alternative trials, improving the robustness of the system. This principle aligns with the notion of self-consistency in LLMs~\cite{wang2022self}, whereby rules recurrent across trials exhibit greater validity. Secondly, the ensemble approach addresses the constraint imposed by the LLM's limited prompt capacity. For scenarios involving tabular datasets surpassing prompt size restrictions, bagging effectively reduces the input scope, thereby preserving dependable performance. The ensemble of multiple weak learners, developed over repeated trials, counterbalances the shortcomings of any single rule set with incomplete feature or instance coverage, ultimately resulting in a comprehensive prediction mechanism.

\section{Experiments}
We assess the effectiveness of \textsf{FeatLLM} across a diverse array of tabular datasets, complemented by several analytical studies aimed at understanding the underlying mechanisms of our framework. Due to limitations on manuscript length, supplementary experimental results—including tests involving alternative LLMs, qualitative assessments, and further analyses—are detailed in the Appendix.

\subsection{Performance Evaluation}
\paragraph{Datasets.}
Our experiments are conducted using 13 tabular datasets spanning binary and multi-class classification scenarios: (1) Adult~\cite{asuncion2007uci}, which aims to predict whether an individual's annual income surpasses \$50,000; (2) Bank~\cite{moro2014data}, targeting the prediction of a customer's likelihood to subscribe to a term deposit; (3) Blood~\cite{yeh2009knowledge}, for forecasting whether a donor will make future donations; (4) Car~\cite{kadra2021well}, which assesses the quality classification of automobiles; (5) Communities~\cite{misc_communities_and_crime_183}, predicting local crime rates; (6) Credit-g~\cite{kadra2021well}, for the evaluation of an individual's creditworthiness (good or bad credit risk); (7) Diabetes\footnote{\texttt{kaggle.com/datasets/uciml/pima-indians-diabetes-database}}, used for diagnosing diabetes; (8) Heart\footnote{\texttt{kaggle.com/datasets/fedesoriano/heart-failure-prediction}}, for predicting the presence of coronary artery disease; and (9) Myocardial~\cite{misc_myocardial_infarction_complications_579}, which focuses on identifying cases of chronic heart failure.

Furthermore, our evaluation incorporates newer datasets and synthetic data that remain relatively unexplored and did not feature in the pre-training data of large language models: (10) Cultivars, which estimates the grain yield category of soybean cultivars\footnote{\texttt{archive.ics.uci.edu/dataset/913}}; (11) NHANES, which predicts the age group corresponding to each subject\footnote{\texttt{archive.ics.uci.edu/dataset/887}}; (12) Sequence-type, designed to classify sequences as arithmetic, geometric, Fibonacci, or Collatz types; (13) Solution-mix, for determining whether the concentration percentage of a four-solution mixture exceeds 0.5. The Cultivars and NHANES datasets were added to the UCI Machine Learning Repository after September 2023, thereby ensuring exclusion from the pre-training corpus for the LLM API (gpt-3.5-turbo-0613) used in our experiments. The Sequence-type and Solution-mix datasets are synthetic, further assuring no prior exposure to the LLM and thus providing a robust testbed for assessment.

These datasets differ in terms of sample counts and inherent complexity, with summary statistics provided in Table~\ref{Tab:basic-desc}. Each dataset is accompanied by well-defined attribute names and descriptions. Datasets such as `Communities' and `Myocardial' present substantial numbers of attributes, often exceeding 100, which can introduce additional challenges due to LLMs’ context length limitations.

\vspace*{-0.12in}
\paragraph{Baselines.}
We benchmark our approach against nine distinct baselines. The initial three represent standard tabular classification methods: (1) Logistic regression (LogReg), (2) XGBoost~\cite{chen2016xgboost}, and (3) RandomForest~\cite{ho1995random}. The fourth and fifth baselines are designed for semi-supervised learning with unlabeled data: (4) SCARF~\cite{bahri2022scarf} and (5) STUNT~\cite{nam2023stunt}. However, the acquisition of extensive unlabeled data can be impractical in many real-world settings. We also include (6) TabPFN~\cite{hollmann2022tabpfn}, which relies on synthetic data with varying distributions to pretrain models. The final three baselines employ LLM-driven strategies: (7) In-context learning (In-context)~\cite{wei2022emergent}, (8) TABLET~\cite{slack2023tablet}, and (9) TabLLM~\cite{hegselmann2023tabllm}. In-context learning entails embedding a handful of training instances directly into the LLM’s input prompt, foregoing parameter adjustments. TABLET enriches inference through external classifiers using rule-sets and prototypes, incorporated as supplementary hints in the prompt. TabLLM adopts parameter-efficient adaptation strategies such as IA3~\cite{liu2022few}, allowing the LLM to be

\vspace*{-0.12in}
\paragraph{Implementation details.}
\textsf{FeatLLM} leverages GPT-3.5 as its large language model, but its framework permits flexibility in LLM selection (refer to Appendix~\ref{sec:backbones} for performance across varied backbones). The LLM inference is conducted with a temperature of 0.5, while the top-p value adheres to the default API setting of 1. We specify ensemble size and rule extraction count as 20 and 10, respectively. The influence of hyper-parameters is summarized in Figure~\ref{fig:hyperparameter2} and further elaborated in Appendix~\ref{sec:hyper-parameter}.
For training the linear model, Adam optimizer is utilized with a learning rate of 0.01, spanning 200 epochs, with $k$-fold cross-validation guiding epoch selection. When reproducing baseline results, GPT-3.5 serves for in-context learning, whereas TabLLM uses T0~\cite{sanh2022multitask} according to its original configuration. Importantly, TabLLM restricts model usage to LLMs with publicly released checkpoints, which precludes GPT-3.5 access. Conventional machine learning models, namely LogReg, XGBoost, and RandomForest, are parameterized via Grid-search, supplemented by $k$-fold cross-validation; $k$ is chosen as either 2 or 4, ensuring representation from all classes within the training split. Additional implementation specifics can be found in Appendix~\ref{sec:baselines}.

\vspace*{-0.12in}
\paragraph{Results.}
Tables~\ref{tab:main1} and \ref{tab:main2} display the comparative performance results across several datasets. Each experiment is conducted three times, using varying random splits of training and test sets, with the reported AUC values including both mean and standard deviation. Across these comparisons, \textsf{FeatLLM} either achieves the highest score or is second-best among all baseline methods. Impressively, our method performs similarly to traditional tabular baseline models that require 32-shot samples, while needing only 4-shot samples (refer to Fig.~\ref{fig:compares}). While models such as STUNT and TabPFN exhibit noticeable gains in certain datasets, their advantages over classic machine learning models are not consistently substantial. Approaches leveraging LLMs—including in-context learning, TABLET, or TabLLM—were unable to effectively perform inference on datasets with large feature sets. Our method, based on bagging and ensemble strategies, maintains solid efficacy even when handling datasets with many features. Additionally, the bagging and ensemble approach we introduce is readily transferable to other LLM-driven frameworks; however, this adaptation would effectively double inference per sample, thus considerably increasing computational demands to a level that is impractical.

\subsection{Analysis \& Discussion}
\label{sec:analysis}
\paragraph{Ablation study.} To discern the contributions of key framework elements, we conduct ablation experiments, investigating the following aspects: (1) \textsf{-Tuning}: removing the linear model’s weight tuning and instead using a sum of binary features for logit computation; (2) \textsf{-Ensemble}: excluding the ensemble component; (3) \textsf{-Description}: omitting feature descriptions; and (4) \textsf{-Reasoning}: skipping Step-1 in reasoning instructions during rule generation.

A summary of how these ablations affect AUC, averaged across datasets, is found in Table~\ref{tab:ablation}. For all modifications, we observe a drop in performance. Notably, the advantage gained from weight tuning amplifies as the number of shots increases, reflecting that abundant data enables more precise rule importance estimation and highlights the impact of tuning. Conversely, when only a few shots are available, the utility of feature descriptions and reasoning instructions rises, implying that maximizing the LLM’s prior knowledge is critical for enhancing outcomes. Finally, the ensemble method delivers consistent improvements across all experimental conditions.

\vspace*{-0.12in}
\paragraph{Training \& inference time comparison.} We evaluate and contrast the training and inference times across various methods. All experiments are conducted on the Adult dataset, utilizing a training set with 16 samples per class (shots). The default configuration leverages a single A100 GPU, with the exception of TabLLM, which employs four A100 GPUs to support model parallelism. Table~\ref{tab:cost} reports the respective runtimes. Of particular interest, \textsf{FeatLLM} demonstrates significantly reduced inference time, remaining in close alignment with that observed for traditional methods. Meanwhile, alternative LLM-based strategies, such as In-context learning, TABLET, and TabLLM, are marked by much greater inference overhead. With regard to training duration, \textsf{FeatLLM} maintains parity with other LLM-based approaches; its requirement of only 30 API queries is not financially prohibitive and enables feasible parallelization.

\vspace*{-0.12in}
\paragraph{Quality of features generated by \textsf{FeatLLM}.}
To assess the practical utility of rule-based features created by \textsf{FeatLLM}, we design a straightforward experiment examining their informativeness for downstream tasks. Specifically, we compute the average AUROC between the newly extracted features and the associated target labels, and determine the proportion of features whose AUROC exceeds 0.5 for each dataset. 
Table~\ref{tab:quality} summarizes these findings. The data verify that most generated rules are substantially relevant to the target variable and offer meaningful contributions to task performance (see the ``Before tuning'' column). Additionally, as the linear model is adapted using the data, features that exhibit minimal connection (i.e., rules with assigned weights falling below a specified threshold) are systematically removed (see the ``After tuning'' column). Here, a threshold value of 0.1 was chosen based on the weight distribution histogram.

\vspace*{-0.12in}
\paragraph{Impact of introducing spurious correlations.} To evaluate how well different approaches eliminate irrelevant spurious correlations under low-shot conditions, we performed a targeted study. In this setup, columns randomly drawn from the Adult dataset—unrelated to the Heart dataset's prediction objective—were appended to the Heart dataset. Model accuracy was then tracked as the quantity of these extraneous columns increased. Ensuring a fair assessment, all models operated with only 8 labeled examples, and did not have access to any unlabeled samples, which precludes consideration of algorithms such as SCARF and STUNT.

Figure~\ref{fig:spurious} presents how model outcomes shift as spurious correlations are introduced. Relative to alternative methods, \textsf{FeatLLM} exhibited the smallest decline in performance, suggesting that its rule extraction process—grounded in leveraging prior knowledge about task-feature relevance—successfully discourages use of features that are not meaningful for the prediction problem. Consequently, rules derived from the noisy, irrelevant columns constituted merely 1-2\% of the total extracted. Nonetheless, \textsf{FeatLLM} may still inadvertently capture and propagate spurious associations—particularly when columns drawn from similar datasets are randomly introduced—leading to confusion between genuine causal connections and incidental patterns (refer to Appendix~\ref{sec:spurious-appendix} for details).

\vspace*{-0.12in}
\paragraph{Hyper-parameter analysis}
Two primary hyper-parameters govern \textsf{FeatLLM}: the ensemble count and the number of rules per class extracted from each LLM interaction. To evaluate how these hyper-parameters influence model effectiveness, we performed systematic experiments. We observed that increasing the ensemble count enhances performance, but after approximately 20 ensembles, the improvements plateau (refer to Fig.~\ref{fig:appendix-hyperparameter1} in Appendix). Regarding the quantity of extracted rules, results did not show significant differences; we concluded that extracting 10 rules represents an optimal compromise between prompt size and accuracy (refer to Fig.~\ref{fig:hyperparameter2}). Our hypothesis is that extracting more rules increases the input dimensionality, which adds informational value but may also exacerbate overfitting under limited data conditions.

\section{Conclusion}
This work introduces \textsf{FeatLLM}, which establishes a new standard for LLM-driven few-shot learning with tabular datasets. Unlike approaches where LLMs produce sample-specific predictions, \textsf{FeatLLM} emphasizes feature engineering by generating predictive rules and employing a bagging-based ensemble. This methodology ensures low inference overhead, eliminates the need for additional training, and functions efficiently with API-only access, regardless of the feature count in the data. As a result, \textsf{FeatLLM} attains strong predictive accuracy and emerges as an effective, resource-conscious solution for deploying LLMs in real-world tabular scenarios.

\textsf{FeatLLM} has been optimized specifically for low-shot learning contexts. Looking ahead, we intend to broaden its utility by generating novel feature types for data with more samples, diversifying approaches beyond rule generation, and enhancing interpretability to facilitate wider applicability across diverse tasks.

\section{Broader Impact} Our proposed approach, \textsf{FeatLLM}, leverages the extensive background knowledge contained within LLMs for tabular tasks, outperforming conventional supervised methods in terms of achievable accuracy and sample efficiency. By drawing on pretrained knowledge, the framework addresses overfitting challenges that arise from small labeled datasets. This is especially valuable in practical domains such as finance and healthcare, where labeling requires substantial effort or is otherwise impractical, and where relevant pretrained information may already be embedded in LLMs. A crucial area for subsequent investigation concerns societal biases and misinformation within LLMs' prior knowledge, as these factors may disproportionately shape prediction outcomes and, at times, override information provided by observed labeled instances.

\end{document}